% ---------------------------------------------------------------------
% v3 section: Certified enclosures for the densities.
%
% Every numeric in this section was computed for this draft by
%   prototype/explore/race_density_certified.py
% (absolute path in this repository:
%   /home/user/hihi/algorithms/prime-sum/prototype/explore/race_density_certified.py;
% python3 + mpmath 1.4.1 at 40 working digits with 100 guard digits inside
% the Hurwitz-zeta representation of L, numpy/scipy for the quadrature;
% total runtime ~14 s).  The script re-verifies at run time: the s=2
% anchors L(2,chi_4)=Catalan and the direct series for L(2,chi_3); the
% L'/L(1) anchors by generalized Stieltjes constants against the
% closed form and against numerical differentiation; sigma^2(0) against
% its closed form; sigma_m^2 against the 12-digit reference values of
% v3/verify_variance_identity.py; the constant c_4 below and its sharp
% value; the Bessel envelope constant; the monotonicity of a_gamma; the
% positivity of every tail variance; and a Riemann--von Mangoldt count of
% both ordinate lists.
%
% REPLACEMENT MAP (line numbers refer to the 2026-08-15 draft of
% weighted-races.tex):
%   * The subsection below is to be \input inside
%     \section{Methods and verification}, REPLACING the two sentences
%     "Densities: Monte Carlo ... against \cite{RS94}." (lines 796-800).
%   * Block R1 at the end of this file REPLACES the clause at lines
%     106-108 ("and every headline number was computed by two
%     independent routes").
%   * Block R2 REPLACES the sentence at lines 814-815 ("Every headline
%     number in this note was re-computed independently of the code that
%     first produced it.").
%   * Table~\ref{tab:enclose} SUPERSEDES the numerical entries of
%     Table~\ref{tab:delta} (which are kept there as the historical
%     six-decimal values; see the "corrected values" paragraph for the
%     ten entries whose sixth decimal moves).
%   * The paragraph "Correction to Remark [rem:tails]" CORRECTS the
%     closing claim of that remark in v3/variance-identity.tex.
% Requires only amsmath/amssymb/amsthm/booktabs and the theorem
% environments and bibliography keys (Dav00, RS94, FM13, ANS14) of the
% main file.  Notation as in Section~\ref{sec:theorem} and
% Section~\ref{sec:variance}: M=m+1/2, a_gamma(m)=2M/sqrt(M^2+gamma^2),
% sigma_m^2 = sum 2 a_gamma^2 = 4M (log Lambda)'(m+1).
% ---------------------------------------------------------------------

\subsection*{Certified enclosures for the densities}
\label{subsec:certified}

The densities of Table~\ref{tab:delta} were produced with a tail model:
the unseen zeros $\gamma>\Gamma$ ($\Gamma=\gamma_{511}=639.9157\ldots$
for $\chi_4$, $\gamma_{537}=699.7747\ldots$ for $\chi_3$) were replaced
by a Gaussian whose variance was \emph{estimated} from the
zero-counting density. This subsection replaces the estimate by a
certified computation and states precisely which error terms are proved
and which remain estimated. Throughout, the standing hypotheses are
GRH$(\chi)$ + LI$(\chi)$ (as for every density in this literature)
together with the completeness of the computed ordinate list up to
$\Gamma$ (validated by the Hardy $Z$-scans and Riemann--von Mangoldt
counts described above; the enclosure script repeats a smoothed count
check, which matches to $0.6$ resp.\ $1.3$ of a zero, within the
$S(T)$ fluctuation).

Under LI, $\delta_q(m)=\mathbb P(1+V>0)$ with
$V=\sum_{\gamma>0}2a_\gamma(m)\cos\theta_\gamma$, whose characteristic
function factors over the computed and unseen ordinates,
\[
\Phi(t)\;=\;\prod_{\gamma>0}J_0\bigl(2a_\gamma t\bigr)
\;=\;P(t)\,\Psi(t),\qquad
P=\prod_{\gamma\le\Gamma},\quad \Psi=\prod_{\gamma>\Gamma}.
\]
Since $|J_0(x)|\le\min\{1,\sqrt{2/(\pi x)}\}$ for $x>0$ (the classical
envelope; the extremal value $\sup_{x>0}\sqrt x\,|J_0(x)|=\sqrt{2/\pi}
=0.79788\ldots$ is re-verified numerically in the script, which uses
the rounded-up constant $0.8$), $\Phi$ is integrable, $V$ has a
continuous density, and the Gil--Pelaez inversion is an identity:
\begin{equation}\label{eq:gilpelaez}
\delta_q(m)\;=\;\frac12+\frac1\pi\int_0^\infty
\Phi(t)\,\frac{\sin t}{t}\,dt .
\end{equation}
The enclosure controls the three gaps between \eqref{eq:gilpelaez} and
what a finite computation can evaluate.

\emph{(i) The tail variance is certified, not estimated.} By
Proposition~\ref{prop:varid},
\[
V_T\;:=\;\sum_{\gamma>\Gamma}2a_\gamma(m)^2
\;=\;4M\,\frac{\Lambda'}{\Lambda}(m+1)-\sum_{\gamma\le\Gamma}2a_\gamma(m)^2 ,
\]
with no reference to the zero-counting density. The subtracted finite
sum inherits the $25$-digit accuracy of the ordinates (effect
$<10^{-20}$); the closed form is evaluated by mpmath at $40$ digits
through two independent routes that must agree to $25$ digits, anchored
at $L(2,\chi_4)=G$ (Catalan), at a directly summed series for
$L(2,\chi_3)$, at $(L'/L)(1,\chi_q)$ expressed through generalized
Stieltjes constants, and at the closed form for $\sigma_0^2$. This is
high-precision floating point, \emph{not} interval arithmetic; we
budget a defensive slack of $10^{-15}$ for it, absorbed into the error
radius below. The script aborts if any $V_T\le0$; in fact
$V_T$ ranges from $0.0030943$ ($q=3$, $m=0$) to $139.71$
($q=4$, $m=100$); see Table~\ref{tab:enclose}.

\emph{(ii) The unseen product is a Gaussian with a proved quartic
correction.} Since $a_\gamma=2M/\sqrt{M^2+\gamma^2}$ is strictly
decreasing in $\gamma$, every unseen ordinate satisfies
$2a_\gamma t\le 2a_\Gamma t\le1$ whenever
$t\le t_{\max}:=1/(2a_\Gamma)$. On $|z|\le1$ we use
\begin{equation}\label{eq:quartic}
0\;\le\;-\log J_0(z)-\frac{z^2}{4}\;\le\;c_4\,z^4,
\qquad
c_4=-\log J_0(1)-\tfrac14=0.0176210647\ldots,
\end{equation}
with the \emph{sharp} constant: from the product
$J_0(z)=\prod_n(1-z^2/j_{0,n}^2)$ over the Bessel zeros,
$-\log J_0(z)=\sum_{k\ge1}(\sigma_k/k)z^{2k}$ with Rayleigh sums
$\sigma_k=\sum_nj_{0,n}^{-2k}>0$, valid for $|z|<j_{0,1}=2.40\ldots$;
$\sigma_1=\tfrac14$ (Rayleigh), and $z^{2k}\le z^4$ for $k\ge2$,
$|z|\le1$, so the middle member of \eqref{eq:quartic} equals
$\sum_{k\ge2}(\sigma_k/k)z^{2k}\in[0,c_4z^4]$ with
$c_4=\sum_{k\ge2}\sigma_k/k$, attained at $z=1$. (The script uses the
rounded-up $c_4'=0.018$ and re-verifies \eqref{eq:quartic} on a grid.)
Summing over the unseen zeros: for $0\le t\le t_{\max}$,
\[
\Psi(t)=\exp\Bigl(-\tfrac12V_Tt^2+\eps(t)\Bigr),
\qquad
-c_4'\,S_4\,t^4\le\eps(t)\le0,
\qquad
S_4:=\!\!\sum_{\gamma>\Gamma}\!(2a_\gamma)^4
\le(2a_\Gamma)^2\cdot 2V_T,
\]
the last bound because $(2a_\gamma)^2\le(2a_\Gamma)^2$ beyond $\Gamma$
and $\sum_{\gamma>\Gamma}(2a_\gamma)^2=2V_T$. The sign information is
exploited: the Gaussian model \emph{over}states $|\Psi|$, and with
$w(t):=c_4'S_4t^4$ the factor $e^{\eps(t)}\in[e^{-w},1]$ is enclosed by
its midpoint,
\[
\Bigl|e^{\eps(t)}-\tfrac{1+e^{-w(t)}}2\Bigr|\;\le\;\tfrac{1-e^{-w(t)}}2,
\]
which halves the tail-model radius relative to the symmetric bound
$e^{w}-1$.

\emph{(iii) Enclosure.} With $G(t):=e^{-V_Tt^2/2}$ and a truncation
point $T_1\le t_{\max}$,
\[
\delta_q(m)\;=\;\frac12+\frac1\pi\Bigl(\,
\underbrace{\int_0^{T_1}\!P\,G\,\frac{1+e^{-w}}2\,\frac{\sin t}{t}\,dt}_{I,\ \text{computed}}
\;+\;R\Bigr),
\qquad |R|\;\le\;E_1+E_2,
\]
\[
E_1=\int_0^{T_1}\!|P|\,G\,\frac{1-e^{-w}}2\,dt
\quad(\text{tail-model radius, from the midpoint enclosure of }e^{\eps}),
\]
\[
E_2=\frac{(0.8)^K}{\sqrt{\prod_{i\le K}2a_{\gamma_i}}}\cdot
\frac{T_1^{-K/2}}{K/2}
\;\ge\;\Bigl|\int_{T_1}^\infty\!\Phi\,\frac{\sin t}{t}\,dt\Bigr|
\quad(\text{truncation; envelope on the first }K\text{ ordinates,}
\]
all other factors bounded by $1$, $|\sin t/t|\le1/t$). The script
chooses $K\in\{24,\dots,72\}$ and $T_1\le t_{\max}$ so that
$E_2\le\pi\times10^{-15}$; this holds at every $(q,m)$ computed.
What remains \emph{estimated} rather than certified, stated openly:
the quadrature error of $I$ and of $E_1$ (composite Simpson on $2^{17}$
subintervals in IEEE double precision; a step-doubling Richardson
comparison is below $10^{-12}$ at every $(q,m)$, and is budgeted as
$E_3=10^{-12}$; $E_1$, itself a bound of size $\lesssim5\times10^{-7}$,
is inflated by $50\%$ against its own quadrature error), the
double-precision rounding of the $511$-fold product ($\lesssim10^{-13}$
relative), and the $10^{-15}$ variance slack of (i), together budgeted
as $E_4=2\times10^{-11}$ plus the computed slack sensitivity. The
radius of the quoted intervals is $(E_1+E_2+E_3+E_4)/\pi$, dominated
everywhere by the proved bound $E_1$. These enclosures are rigorous
\emph{modulo} exactly the items just listed (double-precision
quadrature with Richardson control and the anchored $40$-digit
evaluation of $\sigma_m^2$); no interval arithmetic was used, and we do
not claim otherwise.

\begin{table}[h]\centering
\begin{tabular}{crrllcr}
\toprule
$q$ & $m$ & $V_T$\quad\ & $\delta_q(m)$ lower & $\delta_q(m)$ upper &
width & Tab.~\ref{tab:delta} $-$ mid \\
\midrule
$4$ & $0$   & $0.0034836$ & $0.9959278380$ & $0.9959280127$ & $1.7\times10^{-7}$ & $+0.07$ \\
$4$ & $1$   & $0.0313527$ & $0.7976303685$ & $0.7976304296$ & $6.1\times10^{-8}$ & $-2.40$ \\
$4$ & $2$   & $0.0870906$ & $0.6945210810$ & $0.6945211213$ & $4.0\times10^{-8}$ & $-1.10$ \\
$4$ & $3$   & $0.1706968$ & $0.6458510415$ & $0.6458510746$ & $3.3\times10^{-8}$ & $-1.06$ \\
$4$ & $8$   & $1.0067182$ & $0.5718700860$ & $0.5718701185$ & $3.3\times10^{-8}$ & $-1.10$ \\
$4$ & $20$  & $5.8541890$ & $0.5384552785$ & $0.5384553290$ & $5.1\times10^{-8}$ & $-0.30$ \\
$4$ & $100$ & $139.71245$ & $0.5137782134$ & $0.5137783837$ & $1.7\times10^{-7}$ & $-0.30$ \\
\midrule
$3$ & $0$   & $0.0030943$ & $0.9990630500$ & $0.9990633364$ & $2.9\times10^{-7}$ & $-0.19$ \\
$3$ & $1$   & $0.0278484$ & $0.8368804319$ & $0.8368805308$ & $9.9\times10^{-8}$ & $-3.48$ \\
$3$ & $2$   & $0.0773564$ & $0.7234728791$ & $0.7234729425$ & $6.3\times10^{-8}$ & $-2.91$ \\
$3$ & $3$   & $0.1516181$ & $0.6665542399$ & $0.6665542896$ & $5.0\times10^{-8}$ & $-2.26$ \\
$3$ & $8$   & $0.8942043$ & $0.5785355754$ & $0.5785356134$ & $3.8\times10^{-8}$ & $-1.59$ \\
$3$ & $20$  & $5.2001295$ & $0.5407681867$ & $0.5407682369$ & $5.0\times10^{-8}$ & $-1.21$ \\
$3$ & $60$  & $45.202127$ & $0.5197174290$ & $0.5197175306$ & $1.0\times10^{-7}$ & $-1.48$ \\
\bottomrule
\end{tabular}
\caption{Certified enclosures $[\delta_{\mathrm{lo}},
\delta_{\mathrm{hi}}]$ for the logarithmic densities, under
GRH + LI and completeness of the ordinate lists to height $\Gamma$
($511$ ordinates for $\chi_4$, $537$ for $\chi_3$). $V_T$ is the
certified variance of the unseen-zero tail. The last column is the
deviation of the previously printed six-decimal value from the interval
midpoint, in units of $10^{-6}$; entries exceeding $1$ in absolute
value are discrepancies beyond the sixth decimal (ten of the fourteen).
All widths are below $3\times10^{-7}$, against the design target
$5\times10^{-6}$.}
\label{tab:enclose}
\end{table}

\emph{Corrected values, and what was wrong.} The enclosures confirm
Table~\ref{tab:delta} at $(q,m)=(4,0)$, $(4,20)$, $(4,100)$, $(3,0)$
and move the sixth decimal in the other ten entries by $1$--$4$ units;
the corrected six-decimal values are
\[
\delta_4:\ 0.995928,\ 0.797630,\ 0.694521,\ 0.645851,\ 0.571870,\
0.538455,\ 0.513778;
\]
\[
\delta_3:\ 0.999063,\ 0.836880\,(*),\ 0.723473,\ 0.666554,\ 0.578536,\
0.540768,\ 0.519717\,(*),
\]
$(*)$: two sixth decimals are \emph{not resolved} by the enclosures.
At $(3,1)$ the interval $[0.83688043,0.83688053]$ straddles the
rounding boundary $0.8368805$ (the value is $0.836880$ or $0.836881$;
the midpoint $0.83688048$ sits within $2.5\times10^{-9}$ of the
boundary, so no computation at this zero height can be expected to
settle it). At $(3,60)$ the interval $[0.51971743,0.51971753]$
straddles $0.5197175$ (the value is $0.519717$ or $0.519718$). All
twelve remaining sixth decimals, including $(4,100)$, are certified by
the enclosures. The cause of the drift is
quantitatively identified: re-running the certified pipeline with the
old smoothed-tail variance in place of $V_T$ reproduces the deviations
(e.g.\ $-3.2\times10^{-6}$ of the observed $-3.5\times10^{-6}$ at
$(3,1)$, the small remainder being the old code's coarser quadrature).
The classical anchors are unaffected: our intervals contain
$0.9959279\ldots$ and $0.9990632\ldots$ at $m=0$, consistent with
\cite{RS94,FM13} to all six printed decimals.

\emph{Correction to Remark~\ref{rem:tails}.} That remark asserted that
replacing the estimated tail by the exact one moves the densities of
Table~\ref{tab:delta} by less than their quoted precision. The
enclosures above show this is false at the sixth decimal for ten of
fourteen entries (largest shift $3.5\times10^{-6}$, at $(q,m)=(3,1)$);
the remark's claim should be read as ``below $5\times10^{-6}$'', and
Table~\ref{tab:enclose} now supersedes the affected entries.

\emph{The two computations, honestly described.} The deterministic
Gil--Pelaez inversion and the Monte Carlo evaluation of
Section~\ref{sec:methods} share their inputs: the identical
$511$- resp.\ $537$-ordinate lists and (in their original form) the
identical Gaussianized tail. Their mutual agreement to
$\le10^{-4}$ therefore validated only the \emph{evaluation} step ---
quadrature against sampling --- and could not have detected an error in
the shared tail model. They are independent as implementations, not as
models. The enclosures of this subsection close that gap from the model
side: the tail variance is now the certified $V_T$ of (i), the
non-Gaussianity of the tail is bounded by the proved $E_1$, and a
Monte Carlo rerun with the certified $V_T$ in the tail agrees with the
enclosure midpoints within half a standard error at each tested point
($2\times10^6$ samples at $(q,m)=(4,0),(4,1),(3,1)$).

\emph{General real primitive modulus.} For use beyond $q\in\{3,4\}$ we
record the general form of Theorem~\ref{thm:interp} and of the variance
identity. Let $\chi$ be a real primitive character mod $q$ with
$\chi(-1)=(-1)^a$, $a\in\{0,1\}$, and assume GRH$(\chi)$, LI$(\chi)$,
and the no-central-zero hypothesis
\[
\mathrm{NCZ}(\chi):\qquad L(\tfrac12,\chi)\neq0 .
\]
Define $R_m(x)=-\sum_{p\le x}\chi(p)p^m\log p$ (nonresidues minus
residues), $E_m=(2m+1)x^{-(m+1/2)}R_m$, and
$a_\gamma(m)=2M/\sqrt{M^2+\gamma^2}$ over the positive ordinates of
$L(s,\chi)$. Then parts (i)--(iii) of Theorem~\ref{thm:interp} hold
verbatim with $C(\chi)=\sum_{\gamma_\chi>0}2\gamma_\chi^{-2}$, and the
variance identity becomes, with
$\xi(s,\chi)=(q/\pi)^{(s+a)/2}\Gamma\bigl(\tfrac{s+a}2\bigr)L(s,\chi)$,
\[
\sigma_m^2(\chi)\;=\;4M\,\frac{\xi'}{\xi}(m+1,\chi),
\qquad
C(\chi)\;=\;(\log\xi)''(\tfrac12,\chi).
\]
What changes and what does not. \emph{Unchanged:} the prime-square bias
term --- the $k=2$ prime powers contribute
$\sum_{p\le\sqrt x,\,p\nmid q}p^{2m}\log p=x^{m+1/2}/(2m+1)+o(x^{m+1/2})$
because $\chi(p)^2=1$ for every $p\nmid q$ and the finitely many $p\mid q$
contribute $O_m(x^{m}\log x)$, so the normalized mean is $1$ for every
real $\chi$; the Chernoff/Chebyshev bounds of (iii), which use only the
arcsine moment generating function. \emph{Changed:} the gamma factor
and conductor in $\xi$ (parity $a$ enters the variance asymptotic only
through the constant term $2a-1$ of
Proposition~\ref{prop:varasymp}); all implied constants of the
explicit-formula lemmas, which now depend on $q$, and the
first-ordinate constants ($|\rho+m|\ge\gamma_1(\chi)$ in
Lemma~\ref{lem:trunc}; numerically $\gamma_1(\chi_4)=6.0209\ldots$,
$\gamma_1(\chi_3)=8.0397\ldots$, and $\gamma_1(\chi)>0$ is exactly
NCZ under GRH); for $a=0$ the pole of $1/s$ in the Perron contour meets
the trivial zero of $L(s,\chi)$ at $s=0$, and the classical formula
carries an extra $-\log x+b(\chi)$ term \cite[Ch.~19]{Dav00}, absorbed
by the $O(\log x)$ floor of \eqref{eq:trunc}; and the functional
equation input, $\xi(s,\chi)=\xi(1-s,\chi)$, i.e.\ root number $+1$,
which holds for every real primitive $\chi$ by the evaluation of the
sign of the quadratic Gauss sum ($\tau(\chi)=\sqrt q$ for $a=0$,
$i\sqrt q$ for $a=1$; \cite[Ch.~2]{Dav00} for prime modulus, the
general case reducing to it via the factorization of Gauss sums over
coprime moduli --- for the moduli $3,4$ used here the two-term sums
are evaluated directly in the verification script). On the
no-central-zero hypothesis itself: for $\chi_4$ it is the alternating
series $L(\tfrac12,\chi_4)=\sum_k(-1)^k(2k+1)^{-1/2}>0$; for $\chi_3$
it is \emph{not} an alternating-series statement (the zero terms
$\chi_3(3k)=0$ interleave), but it is still unconditional: the
Dirichlet series $\sum\chi_3(n)n^{-1/2}$ converges by Dirichlet's test
(bounded character sums) and equals $L(\tfrac12,\chi_3)$ by analytic
continuation, and regrouping the convergent series in blocks of three
gives
\[
L(\tfrac12,\chi_3)\;=\;\sum_{k\ge0}
\Bigl[\frac1{\sqrt{3k+1}}-\frac1{\sqrt{3k+2}}\Bigr]
\;\ge\;1-\frac1{\sqrt2}\;>\;0
\]
(every block positive; numerically $0.4808675\ldots$). For general
$q$ we do not prove NCZ: under GRH one has $L(\tfrac12,\chi)\ge0$
(a negative central value would force a real zero in $(\tfrac12,1)$),
so NCZ is then exactly central nonvanishing, which is Chowla's
conjecture for quadratic characters --- known in every numerically
checked case but open in general; it must be assumed, or verified
computationally for the modulus at hand, as part of the hypotheses.
Nothing in this paper uses NCZ beyond $q\in\{3,4\}$, where it is proved
above. For the record, the certified constants at the two moduli used
here are
\[
C(\chi_4)=0.15602964988964488\ldots,\qquad
C(\chi_3)=0.11340215685275857\ldots,
\]
\[
\sigma_0^2(\chi_4)=0.15556797992358593\ldots,\qquad
\sigma_0^2(\chi_3)=0.11322996985747234\ldots
\]

% ---------------------------------------------------------------------
%% Block R1.  REPLACES the clause at lines 106-108 of weighted-races.tex
%  ("..., and every headline number was computed by two independent
%  routes.").
% ---------------------------------------------------------------------
%
% ..., at the standard of numerical work in this area (GRH + LI where
% densities are involved). Every headline number was computed twice: the
% sieve counts by fully independent implementations, the densities by
% deterministic inversion against Monte Carlo --- two evaluations that
% share the same zero data and tail model and are independent only in
% the evaluation step; the shared tail model is bounded rigorously by
% the certified enclosures of Section~\ref{sec:methods}
% (Table~\ref{tab:enclose}).
%
% ---------------------------------------------------------------------
%% Block R2.  REPLACES the sentence at lines 814-815 of
%  weighted-races.tex ("Every headline number in this note was
%  re-computed independently of the code that first produced it.").
% ---------------------------------------------------------------------
%
% Every headline number in this note was re-computed by a second
% implementation. For the sieve computations the two implementations are
% independent throughout; for the densities they shared the ordinate
% lists and, prior to the enclosures above, the Gaussian tail model, so
% their agreement tested evaluation error only --- the tail model itself
% is certified by Table~\ref{tab:enclose}.
